\section{应用与联系}\label{sec:applications}

\subsection{统计推断的逻辑基础}

统计学的根本问题——\textbf{用样本推断总体}——在大数定律中获得第一原理地位：

\begin{itemize}
    \item \textbf{频率 $\leftrightarrow$ 概率}：Glivenko--Cantelli 定理（\cref{thm:gc}）保证经验分布函数一致收敛于总体分布，使“经验分布可作为总分布的代理”成为定理而非信仰；由此出发，凡能写成样本平均的统计量（样本矩、样本分位数、MLE）都自动获得（弱）相合性。
    \item \textbf{极大似然估计的相合性}：MLE 的对数似然是独立同分布随机变量之和，其期望在真参数处被 Kullback--Leibler 散度的非负性严格支配——大数定律保证似然比在真参数处压倒一切错误参数（Wald 1949 的相合性证明的核心机制）。
    \item \textbf{一致大数定律与学习理论}：当函数类 $\mathcal{F}$ 无穷时，逐点的 $\sup$ 与极限不可交换；Vapnik--Chervonenkis 理论证明，$\mathcal{F}$ 的 VC 维有限时一致收敛 $\sup_{f \in \mathcal{F}} |P_n f - Pf| \to 0$ a.s. 成立——这正是经验风险最小化可学习的充要条件。VC 定理因此被称为“一致大数定律”，Glivenko--Cantelli 定理是其 $\mathcal{F}$ 取示性函数类的特例。
\end{itemize}

\subsection{蒙特卡洛方法}

蒙特卡洛积分以 $\hat{I}_n = \frac{1}{n}\sum_{k=1}^n g(U_k)$ 估计 $I = \mathbb{E}[g(U)]$（$U$ 为均匀随机数）：大数定律保证 $\hat{I}_n \to I$，中心极限定理给出误差 $\sqrt{n}(\hat{I}_n - I)$ 趋于正态，从而\textbf{误差以 $O(n^{-1/2})$ 速率下降且与维数无关}。大数定律还支撑着：马尔可夫链蒙特卡洛（MCMC）的遍历平均估计（大数定律的马尔可夫链版本）、重要性抽样与拟蒙特卡洛的误差分析。可以说，现代计算科学中一切“用随机性换确定性”的方法，其正确性证书都由大数定律签发。

\subsection{信息论}

香农的\textbf{渐近均分性}（AEP，Asymptotic Equipartition Property）本质上是信息量 $-\log p(X_k)$ 的大数定律：$-\frac{1}{n}\log p(X_1 \cdots X_n) \to H$ a.s.。由此，典型序列集合以约 $2^{nH}$ 个元素覆盖几乎全部概率质量——这正是信源编码（压缩极限 $H$）与信道编码（容量 $\log M - H$）的定界工具。香农 1948 年的原始证明用的是弱大数定律（切比雪夫不等式），现代处理则借助遍历定理推广到平稳遍历信源（Shannon--McMillan--Breiman 定理）。

\subsection{数论与动力系统}

\begin{itemize}
    \item \textbf{正规数}：博雷尔 1909 年的原始动机之一——几乎每个实数在任意 $b$ 进制下各位数字独立均匀分布，故数字频率的大数定律给出“几乎所有数都是正规数”\cite{borel1909}。对于给定的具体数（如 $\pi$、$\sqrt{2}$）是否正规，大数定律无能为力，这至今是数论未解之谜。
    \item \textbf{等分布与遍历方法}：把序列 $\{x_n\}$ 在 $[0,1)$ 中的分布问题转化为保测变换 $T: x \mapsto x + \alpha$（无理旋转）沿轨道的观测，逐点遍历定理给出 Weyl 等分布定理的遍历证明；连分数展开的度量理论（Gauss 变换的遍历性、Khinchin 常数与 L\'evy 常数）同样是“确定性系统中的大数定律”。
\end{itemize}

\subsection{随机过程、物理与保险}

\begin{itemize}
    \item \textbf{随机游走与马尔可夫链}：常返马尔可夫链的时间平均（访问频率）收敛于平稳分布——马尔可夫 1906 年的原始推广\cite{markov1906}，如今是 MCMC 收敛性理论的基石；更新过程中的平均更新率 $N(t)/t \to 1/\mu$（$\mu$ 为平均更新间隔）是更新定理的大数定律形式。
    \item \textbf{统计物理}：宏观观测量（压强、温度）作为海量微观粒子作用量的平均而稳定，其合法性论证依赖（广义）大数定律；玻尔兹曼与吉布斯的遍历假说正是把“时间平均 = 空间平均”作为物理公设，而伯克霍夫定理\cite{birkhoff1931}给出了它成立的数学条件。
    \item \textbf{精算与风险}：保险业“个体损失不可测、总体损失率可测”的经营原理即大数定律；破产理论（Cram\'er--Lundberg 模型）中净收入过程的长期行为由（重尾情形修正后的）大数定律控制。
\end{itemize}

\subsection{与中心极限定理的分工}

大数定律与中心极限定理（CLT）描述同一列部分和的\textbf{不同尺度}：

\begin{itemize}
    \item 大数定律：$S_n/n \to \mu$，回答“平均是否收敛”——一阶尺度 $n$；
    \item 中心极限定理：$(S_n - n\mu)/\sqrt{n} \Rightarrow N(0, \sigma^2)$，回答“以何种速率、何种分布涨落”——涨落尺度 $\sqrt{n}$；
    \item 迭代对数定律：以 $\sqrt{2n\log\log n}$ 刻画 a.s. 意义下的精确波动上界。
\end{itemize}

三者构成递进的“尺度阶梯”：大数定律给出存在性与定性结论，CLT 给出分布级的定量涨落，迭代对数定律给出轨道级的精确包络。应用中，大数定律负责“可估计性”（相合性），CLT 负责“置信区间与检验”。
